% !TEX root = ../../main.tex
% C4.2 — Các yêu cầu chức năng (VIẾT GỌN 2026-09-29: giữ mã, ràng buộc, số liệu; bản trước ở git)
% Nguồn: project_requirements.md mục 1–3, 5, 6; usecase_detail.md \mbox{UC-01}..\mbox{UC-15} (ghi chú mỗi UC); architect.md mục 2, 4, 6.
% Đã đối chiếu giao diện: sắp xếp/lọc kết quả (features/search/SearchResults.jsx), lọc vụ việc theo trạng thái/ngày tạo/người phụ trách
% (pages/cases/CaseBrowserPage.jsx), tải video dự phòng (pages/admin/VideoProcessingPage.jsx).
% Mã yêu cầu: FR-C (chung), FR-A (Quản trị viên), FR-O (Giám sát viên), FR-V (Quản lý), FR-S (hệ thống tự động). Dùng lại ở Chương 8; bảng 4.5 không ghi mã FR (SV bỏ), đã bỏ mã UC ở cuối mỗi FR (SV yêu cầu 2026-09-28).
\section{Các yêu cầu chức năng}
\label{sec:4.2}

Các yêu cầu chức năng được nhóm theo tác nhân, cùng nhóm chức năng chung và nhóm chức năng hệ thống tự thực hiện. Mã của chúng được dùng trong đặc tả use case (Mục~\ref{sec:4.5}) và kiểm thử (Chương~\ref{chapter:kiemthu}).

\subsection{Chức năng chung}
\label{sec:4.2.1}

\begin{reqlist}
    \item[FR-C01] Người dùng đăng nhập bằng tên đăng nhập và mật khẩu, rồi được đưa tới giao diện của vai trò; tài khoản bị khóa hoặc ngừng hoạt động bị từ chối.
    \item[FR-C02] Người dùng đăng xuất được. Phiên hết hạn sau 12 giờ kể từ khi được cấp hoặc sau 30 phút không hoạt động, mất hiệu lực ngay khi tài khoản bị khóa hoặc ngừng hoạt động, và được thay bằng phiên mới khi người dùng còn thao tác.
\end{reqlist}

\subsection{Chức năng của Quản trị viên}
\label{sec:4.2.2}

\paragraph{Quản lý tài khoản.}
\begin{reqlist}
    \item[FR-A01] Tạo, cập nhật, khóa, mở khóa, xóa hoặc ngừng hoạt động tài khoản, gán vai trò Giám sát viên hoặc Quản lý; tên đăng nhập là duy nhất.
    \item[FR-A02] Mỗi Giám sát viên được gán đúng một khu vực giám sát, chọn từ danh mục có sẵn.
\end{reqlist}

\paragraph{Quản lý camera.}
\begin{reqlist}
    \item[FR-A03] Thêm và cập nhật camera gồm mã, tên, khu vực, địa chỉ RTSP và thông tin xác thực; mã camera là duy nhất, khu vực không đổi được sau khi tạo.
    \item[FR-A04] Kiểm tra kết nối RTSP khi thêm camera, khi đổi địa chỉ và theo yêu cầu; camera không kết nối được vẫn được lưu, kèm cảnh báo.
    \item[FR-A05] Loại camera khỏi vận hành hoặc đưa trở lại. Dữ liệu của camera ngừng vận hành được giữ nhưng không tìm kiếm hay thêm vào vụ việc được, còn kết quả đã lưu trong vụ việc vẫn xem được; camera vận hành trở lại ở trạng thái tắt xử lý AI.
\end{reqlist}

\paragraph{Quản lý xử lý AI.}
\begin{reqlist}
    \item[FR-A06] Bật hoặc tắt xử lý AI cho từng camera; từ chối bật khi camera mất kết nối RTSP, đã ngừng vận hành hoặc mô hình chưa sẵn sàng.
    \item[FR-A07] Tải lên tệp video cho một camera, xử lý qua cùng quy trình như luồng RTSP; dùng cho dữ liệu trình diễn (Mục~\ref{sec:1.3}) và camera không có RTSP.
    \item[FR-A08] Chọn Detector và Tracker cho toàn hệ thống từ danh sách đăng ký sẵn; cặp mô hình chỉ được áp dụng khi cả hai sẵn sàng và tương thích, nếu không thì giữ cấu hình cũ. Bộ mã hóa là cố định.
\end{reqlist}

\paragraph{Theo dõi vận hành.}
\begin{reqlist}
    \item[FR-A09] Theo dõi trạng thái RTSP và xử lý AI của từng camera, lọc các camera bất thường.
    \item[FR-A10] Kiểm tra nhóm xử lý camera, trên khung hình thật của một camera, và nhóm tìm kiếm, với kết quả theo từng thành phần; khung hình không có người được báo là chưa thể xác minh thay vì lỗi.
    \item[FR-A11] Xem và lọc nhật ký hệ thống theo người thực hiện, loại sự kiện và thời gian. Nhật ký ghi đăng nhập, đăng nhập thất bại, đăng xuất, thao tác quản trị, thao tác trên vụ việc và lỗi kỹ thuật quan trọng, nhưng không ghi tìm kiếm.
\end{reqlist}

\subsection{Chức năng của Giám sát viên}
\label{sec:4.2.3}

\paragraph{Tìm kiếm người.}
\begin{reqlist}
    \item[FR-O01] Tìm bằng ảnh, đưa vào bằng chọn tệp, kéo thả hoặc dán; có thể kéo một kết quả vào ô tìm kiếm để tìm tiếp chính người đó.
    \item[FR-O02] Tìm bằng câu mô tả tiếng Anh; không hỗ trợ tiếng Việt và không dịch tự động.
    \item[FR-O03] Tìm theo thuộc tính (giới tính, loại và màu trang phục trên, dưới, vật mang theo, không có lựa chọn phủ định), được ghép thành câu mô tả tiếng Anh rồi tìm như FR-O02.
    \item[FR-O04] Phạm vi tìm kiếm giới hạn trong khu vực của tài khoản, có thể thu hẹp theo camera đang vận hành và khoảng thời gian; yêu cầu ngoài quyền bị từ chối.
    \item[FR-O05] Chọn số kết quả 4, 8, 12 hoặc 16, trả về theo điểm phù hợp giảm dần, không có ngưỡng điểm.
\end{reqlist}

\paragraph{Xem và đánh giá kết quả.}
\begin{reqlist}
    \item[FR-O06] Mỗi kết quả là một lần xuất hiện của một người trên một camera, hiển thị ảnh người trong ngữ cảnh cùng thứ hạng, camera, khu vực, thời điểm và điểm phù hợp.
    \item[FR-O07] Chọn một kết quả để xem khung hình toàn cảnh có khung bao của người và thông tin chi tiết.
    \item[FR-O08] Sắp xếp kết quả theo điểm phù hợp hoặc thời gian, và lọc lại theo camera hoặc thời gian.
\end{reqlist}

\paragraph{Quản lý vụ việc.}
\begin{reqlist}
    \item[FR-O09] Tạo vụ việc mới từ một kết quả với tiêu đề và ghi chú; người phụ trách luôn là tài khoản đang đăng nhập.
    \item[FR-O10] Thêm kết quả vào một vụ việc đang xử lý do mình phụ trách; mỗi lần lưu tạo một mục mới, kể cả khi trùng lần xuất hiện.
    \item[FR-O11] Có thể đánh dấu vụ việc hoàn thành ngay khi lưu; nếu đánh dấu thất bại, kết quả vẫn được lưu.
    \item[FR-O12] Sửa tiêu đề, ghi chú và loại kết quả khỏi vụ việc đang xử lý.
    \item[FR-O13] Đánh dấu vụ việc hoàn thành hoặc mở lại; vụ việc đã hoàn thành không sửa được nội dung và không thêm, loại được kết quả.
    \item[FR-O14] Xem các vụ việc mình phụ trách, kể cả sau khi đổi khu vực; kết quả hiển thị như FR-O06, FR-O07, không có điểm phù hợp.
\end{reqlist}

\subsection{Chức năng của Quản lý}
\label{sec:4.2.4}

\begin{reqlist}
    \item[FR-V01] Xem tổng quan toàn hệ thống: số vụ việc theo trạng thái, số kết quả đã lưu và danh sách vụ việc gần đây; chọn một vụ việc để xem chi tiết.
    \item[FR-V02] Xem mọi vụ việc, lọc theo trạng thái, thời gian tạo hoặc Giám sát viên phụ trách.
    \item[FR-V03] Xem chi tiết một vụ việc và các kết quả đã lưu ở chế độ chỉ đọc.
    \item[FR-V04] Xem chi tiết một kết quả đã lưu và khung hình toàn cảnh có khung bao; không hiển thị điểm phù hợp.
\end{reqlist}

\subsection{Chức năng hệ thống tự động thực hiện}
\label{sec:4.2.5}

\begin{reqlist}
    \item[FR-S01] Phân tích lần lượt mọi camera đang vận hành, có RTSP và đang bật xử lý AI.
    \item[FR-S02] Với mỗi luồng RTSP hoặc tệp video: lấy mẫu khung hình, phát hiện người, theo vết thành các lần xuất hiện, chọn một khung hình đại diện và tạo vector đặc trưng.
    \item[FR-S03] Lưu khung hình đại diện, khung bao, vector, camera, khu vực và thời điểm của mỗi lần xuất hiện; lần xuất hiện chỉ tìm được khi đã lưu đủ (NFR-06).
\end{reqlist}
